% ============================================================
% AGGIUNTE ALLA CLASSE: spiegazioni, diagrammi, marcatori
% ============================================================
%
% UTILIZZO:
%   \begin{spiegazione}[Titolo opzionale] ... \end{spiegazione}
%       riquadro verde che segue ogni dimostrazione: l'idea in parole semplici,
%       con un diagramma o un esempio passo passo (non si spezza tra due pagine)
%   \begin{passi} \item ... \item ... \end{passi}
%       elenco di passi numerati con un cerchietto (per gli esempi passo passo)
%   \senzadim
%       targhetta «senza dimostrazione»: il programma non chiede la dimostrazione
%   \facoltativa
%       targhetta per le dimostrazioni che il programma non chiede ma che aiutano a capire
%
%   Stili TikZ:
%     Hasse:      [hasse] sull'ambiente; nodi \node[hv] (etichette) o \node[hpunto,label=...] (pallini);
%                 linee \draw[hlato]; evidenziazioni hv evid / hv evid2
%     Relazioni:  nodi \node[rv]; frecce \draw[rfreccia]; cappi \draw[rcappio] (n) to (n)
%     Funzioni:   insiemi \node[insieme, fit=...]; elementi \node[elem]; frecce \draw[fun]
% ============================================================

\RequirePackage{tikz-cd}

% Indice: i numeri di pagina a tre cifre non entrano nella casella standard di 1,55em e la
% classe (\overfullrule) segnerebbe ogni riga con un quadratino nero; la casella si allarga
% (a una misura fissa, così tutti i numeri finiscono sullo stesso margine destro).
\makeatletter
\renewcommand{\@pnumwidth}{22pt}
\makeatother
% Nelle righe delle parti la casella del numero, composta in corpo 14, sporgerebbe oltre il
% bordo del foglio: il numero si allinea al margine destro senza casella.
\titlecontents{part}
    [0pt]
    {\addvspace{1cm}\centering}
    {\color{primary!60!black}\fontsize{16}{0}\selectfont\Gobold Parte \thecontentslabel\quad}
    {\color{primary!60!black}\fontsize{16}{0}\selectfont\Gobold}
    {\hfill\color{primary!80!black}\bfseries\Scientifica\fontsize{14}{0}\selectfont\thecontentspage}
    [\addvspace{0.3cm}\hrule height 1.5pt\addvspace{0.5cm}]

% ── Colori per le classi di equivalenza e le parti di una partizione ──
\colorlet{classeA}{primary!16}
\colorlet{classeB}{orange!22}
\colorlet{classeC}{seagreen!22}
\colorlet{classeD}{pink!30}
\colorlet{classeE}{naplesyellow!45}
\colorlet{classeF}{violet!12}

% ── Riquadro SPIEGAZIONE ──────────────────────────────────────
\newtcolorbox{spiegazionebox}[1]{
  enhanced, breakable=false,
  colback=seagreen!5!white, colframe=seagreen!65!black,
  boxrule=0pt, leftrule=3pt, arc=0pt, outer arc=0pt,
  left=3mm, right=2mm, top=1.2mm, bottom=1.4mm,
  before skip=0mm, after skip=6mm,
  fontupper=\normalsize,
}
\NewDocumentEnvironment{spiegazione}{O{Spiegazione}}{%
  \begin{concetto}\vspace{-3mm}%
  \begin{spiegazionebox}{#1}%
  {\color{seagreen!55!black}{\mdseries\iidea}\ \bfseries #1}\par\nopagebreak\vspace{0.6mm}%
  \setlength{\parskip}{0.6em}%
}{%
  \end{spiegazionebox}\end{concetto}%
}

% ── Passi numerati ────────────────────────────────────────────
\NewDocumentCommand{\numcerchio}{m}{\tikz[baseline=(n.base)]\node[circle, fill=seagreen!65!black, text=white,
  font=\scriptsize\bfseries, inner sep=0pt, minimum size=4mm] (n) {#1};}
\newenvironment{passi}{%
  \begin{enumerate}[label=\numcerchio{\arabic*}, leftmargin=7.5mm, labelsep=2mm,
                    itemsep=0.8mm, topsep=0.6mm, parsep=0.3em]%
}{\end{enumerate}}

% ── Targhette ─────────────────────────────────────────────────
\newcommand{\targhetta}[2]{%
  \tikz[baseline=(t.base)]\node[rounded corners=2pt, fill=#1!12, draw=#1!60!black, line width=0.4pt,
    font=\scriptsize\bfseries, text=#1!60!black, inner xsep=3pt, inner ysep=1.5pt] (t) {#2};}
\newcommand{\senzadim}{\targhetta{neutral}{senza dimostrazione}\par\smallskip}
\newcommand{\facoltativa}{\targhetta{neutral}{dimostrazione non richiesta dal programma}\par\smallskip}

% ── Blocchi colorati (classi di equivalenza, parti di una partizione) ──
% \blocco[larghezza]{colore}{percorso}: «macchia» colorata attorno a un gruppo di elementi.
% Il percorso unisce i centri dei nodi:
%   due elementi:  \blocco{classeA}{(n1.center) -- (n2.center)}
%   tre elementi:  \blocco{classeB}{(n1.center) -- (n2.center) -- (n3.center) -- cycle}
%   un elemento:   \blocco{classeC}{(n1.center) -- ++(0.001,0)}
\newcommand{\blocco}[3][6.4mm]{%
  \begin{scope}[on background layer]
    \draw[line width=#1+0.8pt, line cap=round, line join=round, draw=#2!55!black!45] #3;
    \draw[line width=#1, line cap=round, line join=round, draw=#2, fill=#2] #3;
  \end{scope}}

% ── Stili TikZ ────────────────────────────────────────────────
\tikzset{
  % diagrammi di Hasse
  hasse/.style={every node/.append style={font=\small}},
  hv/.style={draw=primary!70!black, fill=primary!7, rounded corners=2pt, inner sep=2pt,
             minimum height=4.6mm, minimum width=4.6mm},
  hv evid/.style={hv, fill=orange!25, draw=orange!85!black},
  hv evid2/.style={hv, fill=seagreen!25, draw=seagreen!75!black},
  hpunto/.style={circle, fill=primary!80!black, inner sep=0pt, minimum size=5pt},
  hlato/.style={draw=black!60, line width=0.75pt},
  hlato evid/.style={draw=orange!85!black, line width=1.4pt},
  % grafi di relazioni binarie
  rv/.style={circle, draw=primary!75!black, fill=primary!7, inner sep=1pt, minimum size=6mm, font=\small},
  rfreccia/.style={-{Stealth[length=2.3mm]}, line width=0.7pt, draw=black!70},
  % cappi: \draw[rfreccia] (n) to[rcappio] (n);
  rcappio/.style={out=120, in=60, min distance=8mm},
  rcappio sotto/.style={out=-120, in=-60, min distance=8mm},
  rcappio sx/.style={out=150, in=210, min distance=8mm},
  rcappio dx/.style={out=30, in=-30, min distance=8mm},
  % diagrammi a frecce delle funzioni
  insieme/.style={draw=black!60, rounded corners=9pt, line width=0.6pt, inner sep=6pt},
  elem/.style={font=\small, inner sep=1.5pt},
  fun/.style={-{Stealth[length=2.2mm]}, line width=0.7pt, draw=primary!80!black},
  fun2/.style={-{Stealth[length=2.2mm]}, line width=0.7pt, draw=cherryred!85!black},
  fun3/.style={-{Stealth[length=2.2mm]}, line width=0.7pt, draw=seagreen!70!black},
}
